\chapter{Binary Classification}

\section{Confusion Matrix}

A confusion matrix classifies a population of size $n$ into two 
columns, with size P and N, described below.  The classification
of the total population is binary, either as P 
or N, thus P~+~N~=~$n$.
Table~\ref{tab:breast_cancer} shows
the 2-by-2 confusion matrix with the P and N columns. 
\begin{itemize}
\item The first column is the subset of the population $n$
that is {\bf known positive}, with size P 
for a condition $C$.  
\item The second column is the remaining subset of 
the population that is {\bf known negative}, with
size N for a condition $C$.  Negative status for
a condition is denoted
with the over set line, $\overline{C}$, 
to indicate negation of
condition $C$.  
\item Note P $\leq n$ and N $\leq n$.  Typically P $< n$
and N $< n$.  Cases of P~=~$n$ or N~=~$n$ are degenerate 
special cases not considered here.  
\end{itemize}
Thus
P = size($C$),
N = size($\overline{C}$),
$n$ = size($C \;\union\; \overline{C}$), and
$0$ = size($C \;\intersection\; \overline{C}$) since
$\emptyset = C \;\intersection\; \overline{C}$, the 
binary status of condition $C$.
\begin{itemize}
\item Given a known positive P for a condition $C$, the 
P group is divided into two subgroups of 
true positives TP and 
false negatives FN, depending on whether 
or not the
test outcome $T$ correctly predicted the truth 
for the positives in the population.  P = TP + FN.
\item Given a known negative N for condition $C$, thus
$\overline{C}$, the 
N group is divided into two subgroups of 
false positives FP and
true negatives TN, depending on whether or 
not the test outcome $T$ correctly predicted
the truth for the negatives in the population. N = FP + TN.
\item A positive test outcome is denoted $T$.  A
negative test outcome is denoted $\overline{T}$.
\end{itemize}

\begin{table}[htb]
\caption{The 2-by-2 confusion matrix for the condition $C$ 
and the result of the test to identify $C$.  Note: the numbers
in the table are for the size of each population subset for
the worked numerical example.}
\label{tab:breast_cancer}
\centering
\begin{tabular}{|c|c|c|} \hline
Total Population & 
Known Condition $C$&
Known Non-Condition $\overline{C}$ \\ 
$n$ = 10,000 & 
P = TP + FN & 
N = FP + TN \\ \hline \hline
% %%
& 
\colorbox{greenfaded}{\bf True Positive} & 
\colorbox{yellowfaded}{\bf False Positive} \\
Test Positive $T$ & TP = 80 & FP = 495 \\
& ``correct affirmation'' 
& ``false alarm'' (Type I error) \\ \hline
% %%
& 
\colorbox{horange} {\bf False Negative} & 
\colorbox{redfaded} {\bf True Negative} \\
Test Negative $\overline{T}$ & FN = 20 & TN = 9,405 \\ 
& ``miss'' (Type II error) & ``correct rejection'' \\ \hline
\end{tabular}
\end{table}

\subsection{Probabilities for the Confusion Matrix}
\begin{itemize}
\item $P(C)$ is the pre-test probability of having condition $C$, 
also known as the {\bf prevalence}.
\item $P(\overline{C})$ is then the pre-test probability of not 
having condition $C$, thus $\overline{C}$.
$P(\overline{C}) = 1 - P(C)$.
\end{itemize} 

%\vspace{0.5cm}
When the pre-test probability is 
combined with a test, the following 
four outcomes are possible:
\begin{itemize}
\item $P(T|C)$ is the probability of a positive test given a 
known condition.  This is a 
\colorbox{greenfaded}{\bf True Positive} color coded in green.
\item $P(T|\overline{C})$ is the probability of a positive test 
given a known non-condition.  
This is a 
\colorbox{yellowfaded}{\bf False Positive} color coded 
in yellow as a warning
that perhaps another test should be conducted that 
might be able to
rule out the false positive result.
\item $P(\overline{T}|C)$ is the probability of a negative test 
given a known condition.  This is a 
\colorbox{horange}{\bf False Negative} color coded in orange.
\item $P(\overline{T}|\overline{C})$ is the probability of a 
negative test given a known non-condition.   
This is a 
\colorbox{redfaded}{\bf True Negative} color coded in pink.
\end{itemize}

\begin{Hexample}{
Suppose a 44-year-old female has a mammogram test. 
The radiologist interprets the test result 
as ``suspicious for malignancy.'' 
Assume that the overall probability for any 44-year-old
female, regardless of a mammogram test result, 
to have breast cancer is 1-in-100, 
or 1\% (breast cancer prevalence). 
Assume the mammogram test is 80\% sensitive 
and 95\% specific, from data
collected from a population of 10,000 females.  
Table~\ref{tab:breast_cancer} shows
the 2-by-2 contingency table for these data.
}
\end{Hexample}

It is often useful to show each of the four above-described
categories a square, where the area is equal 
to the population size, 
as shown in Fig.~\ref{fig:confusion_area}.

\begin{figure}[htb]
  \begin{center}
    % \begin{overpic}[width=0.7\textwidth, grid, tics=10]{fig/confusion_area.pdf}
    \begin{overpic}[width=0.7\textwidth]{fig/confusion_area.pdf}
      \put( 6.5, 81){TP}
      \put(19, 86.3){FP}
      \put( 9, 76.1){\tiny FN}
      \put(48, 41){TN}
    \end{overpic}
  \end{center}
  \caption{Population of true positives (TP=80), false positives (FP=495), 
  false negatives (FN=20), and true negatives (TN=9,405) represented by the area
  of each square.}
  \label{fig:confusion_area}
\end{figure}

%\vspace{0.5cm}
\newpage
Figure~\ref{fig:decision_tree_mammogram} shows the 
population sorted into a binary tree.

% tikz
% http://www.texample.net/tikz/examples/merge-sort-recursion-tree/
% http://www.texample.net/tikz/examples/probability-tree/

% set overall layout of tree
\tikzstyle{level 1}=[level distance=3.0cm, sibling distance=5.0cm]
\tikzstyle{level 2}=[level distance=4.0cm, sibling distance=2.5cm]

% define styles for bags and leaves
\tikzstyle{bag} = [text width=4em, text centered]
\tikzstyle{end} = [circle, minimum width=3pt, inner sep=0pt]

\begin{figure}[htb]
\centering
\begin{tikzpicture}[grow=right, sloped]
%\begin{tikzpicture}[grow=right]
\node[circle, draw] {\tiny $10,000$}
	child {
		node[circle, draw] {\tiny 9,900}
			child {
				node[circle, draw, fill=redfaded,
				     label=right:
				{\scriptsize $P(\overline{C} \cap \overline{T}) = 
				P(\overline{C}) \cdot P(\overline{T}|\overline{C}) = 
				\frac{9,900}{10,000} \cdot
				\frac{9,405}{9,900} = 
				\frac{9,405}{10,000}$}
				    ] {\tiny 9,405}
				edge from parent
				node[above] {\scriptsize $P(\overline{T}|\overline{C}) = 0.95$}
				node[below] {\scriptsize specificity}
			}
			child {
				node[circle, draw, fill=yellowfaded,
				     label=right:
				{\scriptsize $P(\overline{C} \cap T) = 
				P(\overline{C}) \cdot P(T|\overline{C}) = 
				\frac{9,900}{10,000} \cdot
				\frac{495}{9,900} = 
				\frac{495}{10,000}$}
			     	    ] {\tiny 495}
				edge from parent
				node[above] {\scriptsize $P(T|\overline{C}) = 0.05$}
				node[below] {\hspace{1cm}\scriptsize false positive}
			}
		edge from parent
		node[above] {\scriptsize $P(\overline{C})=0.99$}
	}
	child {
		node[circle, draw] {\tiny 100}
			child {
				node[circle, draw, fill=horange,
				     label=right:
				{\scriptsize $P(C \cap \overline{T}) = P(C) \cdot P(\overline{T}|C) = 
				\frac{100}{10,000}\cdot\frac{20}{100} = \frac{20}{10,000}$}
				    ] {\tiny 20}
				edge from parent
				node[above] {\scriptsize $P(\overline{T}|C) = 0.20$}
				node[below] {\scriptsize false negative}
			}
			child {
				node[circle, draw, fill=greenfaded,
				     label=right:
				{\scriptsize $P(C \cap T) = P(C) \cdot P(T|C) = \frac{100}{10,000}
				\cdot\frac{80}{100} = \frac{80}{10,000}$}
			     	    ] {\tiny 80}
				edge from parent
				node[above] {\scriptsize $P(T|C) = 0.80$}
				node[below] {\hspace{1cm}\scriptsize sensitivity}
			}
		edge from parent
		node[above] {\scriptsize $P(C)=0.01$}
		node[below] {\scriptsize prevalence}
	}; % remember trailing semi-colon after last child
\end{tikzpicture}
\caption{Decision tree for mammogram as a test $T$ for breast 
cancer $C$.}
\label{fig:decision_tree_mammogram}
\end{figure}

\newpage
	
Now, what is often of interest is the {\bf post-test probability}.  
For a given person
receiving the test, $T$, what is the probability of having the condition, $P(C|T)$?
In fact, all variations can be calculated:
\begin{itemize}
\item $P(C|T)$ is the probability of a positive condition given a positive test.
\item $P(\overline{C}|T)$ is the probability of a negative condition given a positive test.
\item $P(C|\overline{T})$ is the probability of a positive condition given a negative test.
\item $P(\overline{C}|\overline{T})$ is the probability of a negative condition given a negative test.
\end{itemize}

\vspace{1.5cm}

Now, the probability of being positive after testing positive is equal to the probability of
being positive and testing positive divided by the probability of testing positive
(regardless of being positive or negative for the condition).
The probability of testing positive equals the probability of being positive and 
testing positive {\em plus} 
the probability of being negative and testing positive.  This is
written as
\be
 P(C|T) = \dst\frac{P(C \cap T)}{P(C \cap T) + P(\overline{C} \cap T)} = 
 \frac{\frac{80}{10,000}}{\frac{80}{10,000} + \frac{495}{10,000}} = 
 0.139, 
 \label{eq:ppv} 
\ee
which seems low!  
A patient, upon hearing this result, 
that there is a 13.9\% probability that 
she has breast 
cancer given that she tested positive, does not reassure.  
The test says she is positive, but the probability that 
she actually {\em is} positive
given the positive test seems low.  How can this be?
The large number of false positives, 
$P(\overline{C} \cap T)$, in the denominator
reduces the result for $P(C|T)$. 
The number of true positives, $P(C \cap T)$, is small 
compared to the number of false
positives.  The result in Eq.~(\ref{eq:ppv}) is 
also referred to as the test's
{\bf Positive Predictive Value (PPV)} 
or {\bf Precision}, described in Section~xxx (to come).

Next, consider the remaining combinations.  
The probability of being negative after 
testing positive, by reasoning similar to that 
given above, is written as
\be
\ P(\overline{C}|T) = \dst\frac{P(\overline{C} \cap T)}{P(\overline{C} \cap T) + P(C \cap T)} = 
 \frac{\frac{495}{10,000}}{\frac{495}{10,000} + \frac{80}{10,000}} = 
 0.861 = 1 - P(C|T) \mbox{\ding{52}}.
 \label{eq:false_discovery_rate}
\ee
Again, we see the effect of false positives, $P(\overline{C} \cap T)$.  
The probability of being negative even though
the test is positive is 86.1\%, giving little comfort to the patient who now thinks, 
because of the test results, that
she is positive even though she is negative. 
The result in Eq.~(\ref{eq:false_discovery_rate}) is also 
referred to as the 
test's {\bf False Discovery Rate (FDR)}.  

The probability of being positive after testing negative is 
\be
 P(C|\overline{T}) = \dst\frac{P(C \cap \overline{T})}{P(C \cap \overline{T}) + 
 P(\overline{C} \cap \overline{T})} = 
 \frac{\frac{20}{10,000}}{\frac{20}{10,000} + \frac{9,405}{10,000}} = 
 0.002 
 \label{eq:false_omission_rate}
\ee
This result is reassuring.  A patient with a negative test has a 0.2\% chance of 
actually being positive.  This occurs because the number of true negatives, 
$P(\overline{C} \cap \overline{T})$, in the denominator, saturates any false negatives, 
$P(C \cap \overline{T})$, and drives the result toward zero.  The result 
in Eq.~(\ref{eq:false_omission_rate}) is also 
referred to as the test's
{\bf False Omission Rate (FOR)}.

The probability of being negative after testing negative is easily calculated as 
1 minus the above result.  For explicitness, however, we show
\be
 P(\overline{C}|\overline{T}) = 
 \dst\frac{P(\overline{C} \cap \overline{T})}{P(\overline{C} \cap \overline{T}) + 
 P(C \cap \overline{T})} = 
 \frac{\frac{9,405}{10,000}}{\frac{9,405}{10,000} + \frac{20}{10,000}} = 
 0.998 = 1 - P(C|\overline{T}) \mbox{\ding{52}}, 
 \label{eq:npv}
\ee
which is also reassuring.  Again, the strength of the true negatives, 
$P(\overline{C} \cap \overline{T})$,  overwhelms the 
effect of the false negatives, $P(C \cap \overline{T})$, in the result.
The result in Eq.~(\ref{eq:npv}) is also referred 
to as the test's 
{\bf Negative Predictive Value (NPV)}.  

